\documentclass[../main.tex]{subfiles}

\begin{document}

\section{Conclusion}

We present the first continuous, statewide map of thaw mode across Alaska's permafrost domain. By combining $19,288$ expert-labeled observations with 70 environmental indicators, we produce a prior-free log-evidence index showing whether points on the landscape more closely resemble abrupt or non-abrupt thaw, paired with an explicit area-of-applicability (AoA) bound. Within that area, which covers roughly $90\%$ of the permafrost domain, we find $24.7\%$ of the landscape has features more consistent with abrupt thaw than non-abrupt thaw. This fraction describes present predisposition, not the observed extent, probability, or future occurrence of abrupt thaw. The classifier achieves an AUC-PR of $0.852$, approximately fifteen times greater than the chance floor, and retains an AUC-PR of $0.54$ even when tested on observations at a median of $251$ km away from the nearest training point.

Topography and temperature provide the leading distinction between thaw modes, but their influence varies spatially across Alaska. Flat, low-lying wetlands and steeper uplands both favor abrupt thaw, consistent with distinct pathways associated with common landforms like thermokarst lakes and retrogressive thaw slumps, respectively. No single indicator controls this signal statewide, but baseline temperature is the dominant influence on the model in $45\%$ of in-AoA grid cells, with the rank importance of leading indicators shifting considerably in space as local environmental factors vary.

These local differences aggregate into a sharply partitioned statewide pattern. Roughly two-thirds of the Brooks Range and Seward Peninsula favor abrupt thaw, as does just under half of the Arctic Foothills, while non-abrupt thaw is favored across the forested interior and mountain regions. Thaw mode nevertheless varies widely within each region, retaining landscape-scale structure that physiographic categories and previously mapped thermokarst occurrence potential do not capture. Our map of thaw mode therefore resolves a largely distinct and complementary dimension of permafrost change. This work helps sharpen our estimates of thaw-related risks to Alaskan communities and where the permafrost-carbon feedback will draw most heavily on abrupt thaw under future climate change.

\end{document}
